Two locked operators compute mixed-spot TME composition and decide when the tumor column is reportable. (A)~Protocol path from a TME reference through specificity weights and a platform self-gate to the reportability gate, which writes ABSTAIN or a simplex composition. (B)~Locked operator identities drawn on the panel: gene specificity \(w_g=1-H(q_g)/\log K\), the weighted working signature \(S=e^{g\hat\gamma}(P\odot w)\), and simplex composition \(\hat\pi=\mathrm{simplex}(W^\star)\). (C)~KEEP/ABSTAIN gate: a malignant--neighbor cosine at or above the frozen cutoff \(c^\star=0.80\) returns ABSTAIN and renormalizes the remaining types; a cosine below the cutoff returns KEEP and reports a tumor fraction. (D)~Patient\_ID donor split: mixed spots and the reference are built from disjoint patients, so \(\{\text{spot patient}\}\cap\{\text{reference patient}\}=\emptyset\). (E)~Evaluation metrics recorded for a computed edge: overall RMSE, type-level Pearson correlation, spot-level Pearson correlation, tumor RMSE, and Jensen--Shannon divergence, with a \(B=1000\) interval on the \(t=0\) tumor RMSE. (F)~Six named operator steps in order: signature extraction, specificity weights, the weighted fit, the platform self-gate, the reportability gate, and the simplex close.
